%!TEX root = main.tex

%FIRRTL represents the standardized elaborated circuit that the Chisel HDL produces.
%FIRRTL represents the circuit immediately after Chisel's elaboration but before any circuit simplification. It is designed to resemble the Chisel HDL after all meta-programming has executed. Thus, a user program that makes little use of meta-programming facilities should look almost identical to the generated FIRRTL.

%For this reason,

%In this section, we give an informal introduction to FIRRTL %intermediate representation
%and its lowering transformations.

\subsection{FIRRTL IR}

%\zhilin{Can we find a running example ?}

%Figure~\ref{fig:example} contains an example FIRRTL program.

\begin{figure}[t]
\centering
\begin{tabular}{c}
%[basicstyle=\footnotesize\ttfamily]
\begin{lstlisting}
circuit A :
  module A :
    input clk : Clock
    input rst : UInt<1>
    output io : {flip enable : UInt<1>, calc : UInt<3>}
    clk is invalid
    rst is invalid
    io is invalid
    inst x of B
    x.io is invalid
    x.clk <= clk
    x.rst <= rst
    x.io.enable <= io.enable
    io.calc <= x.io.calc

  module B :
    input clk : Clock
    input rst : UInt<1>
    output io : {flip enable : UInt<1>, calc : UInt<3>}
    clk is invalid
    rst is invalid
    io is invalid
    reg counter : UInt<3>, clock with : (reset => (rst, UInt<3>("h00")))
    when io.enable
      node _T_7 = add(counter, UInt<1>("h01"))
      node _T_8 = tail(_T_7, 1)
      counter <= _T_8
      skip
    io.calc <= counter
\end{lstlisting}
\end{tabular}\vspace{-2mm}
\caption{An example FIRRTL program \label{fig:example} \tl{need to change this}}
\vspace{-2mm}
\end{figure}

A FIRRTL \emph{circuit} comprises a list of modules with a designated top-level module that has the same identifier as the circuit.
Each module represents a hardware block that can
be instantiated one or more times.
%
For instance, the FIRRTL circuit \prettyll{A} shown in Figure~\ref{fig:example} has two modules \prettyll{A} and \prettyll{B},
where the top-level module of the circuit is the module \prettyll{A}.

A \emph{module} consists of a set of ports and a sequence of statements.
Each port has an identifier, a direction (either \prettyll{input} or \prettyll{output}), and a data type.
For instance, the module \prettyll{A} has two input ports \prettyll{clk} and \prettyll{rst} and one output port \prettyll{io}, with
types \prettyll{Clock}, \prettyll{UInt<1>} and  \prettyll{\{flip enable: UInt<1>, calc: UInt<3>\}}, respectively.
Each \emph{statement} defines a wire, register, node, connect (\prettyll{<=}), partial connect (\prettyll{<-}), conditional (\prettyll{when}), or instance (\prettyll{inst of}), and so on.
A \emph{type} for specifying the structure of the data can be either a \emph{ground type,} e.g. \prettyll{Clock} (carrying a clock signal) and \prettyll{UInt<1>} (unsigned integer of bit width 1);
or an \emph{aggregate type}, which is either a \emph{vector type,} e.g. \prettyll{UInt[3]} (a three-element vector of unsigned integers),
or a \emph{bundle type,} e.g. \prettyll{\{flip enable: UInt<1>, calc: UInt<3>\}} comprising two signals of the ground types \prettyll{UInt<1>} and \prettyll{UInt<3>} respectively.


A circuit \emph{component} is a declared element in a FIRRTL circuit that can be referred to by its name or identifier.
Therefore, a component can be a port, node, wire, register, memory, or instance.
A declared component can be connected to or from except for input ports that can only be connected from.
%\zhilin{component: a double check is needed}
%\david{Double checked: ``component'' is not directly declared in the FIRRTL specification,
%but it mostly speaks of ``declared component'', so it needs to be declared;
%components can be connected to (or from, in the case of module input ports and memory read ports).
%The only kind of component that is missing from the above list is the memory.}

\begin{comment}
A FIRRTL \emph{circuit} comprises modules, each of which include several ports and statements. A circuit must specify the name of the (unique) top-level module.
For instance, \fu{in the FIRRTL program shown in Figure~\ref{fig:example}, the top-level module of the circuit is module {\tt A}.}
Each \emph{port} has an identifier, \fu{an direction either \prettyll{input} or \prettyll{output}, and a data type.
For instance, the module {\tt A} has two input ports  {\tt clk} and {\tt rst} and one output port {\tt io}.}
A \emph{module} includes \fu{a list of ports} and  a sequence of statements.
%A \emph{port}  is either an input port or an output port, and has an identifier as well as a type.
Each \emph{statement} defines a wire, register, node, connect (<=), partial connect (<-), conditional (when), or instance (inst of), and so on.
%A \emph{statement} is either a declaration of wire, register, instance, node, or a connection (<=) or partial connection (<-), or a conditional (when) statement.
A \emph{type} can be either a \emph{ground type}, e.g. \begin{lstinline}!UInt<1>!\end{lstinline} (unsigned integer of bit width 1), or an \emph{aggregate type}, which is either a \emph{vector type}, e.g. \begin{lstinline}!UInt[3]!\end{lstinline} (a three element vector of unsigned integers), or a \emph{bundle type} e.g \begin{lstinline}! {flip enable: UInt<1>, calc: UInt<3>} !\end{lstinline} which comprises two signals of the ground types  \begin{lstinline}!UInt<1>!\end{lstinline} and  \begin{lstinline}!UInt<3>!\end{lstinline} respectively. A circuit \emph{component} is an element in a FIRRTL circuit that can be referred to by its name or identifier. Therefore, a component can be a port, node, wire, register, or instance, and so on. \zhilin{component: a double check is needed}
\end{comment}

The syntactic rules of FIRRTL are given in Figure~\ref{fig-firrtl-synt-simp} in Extended Backus-Naur Form (EBNF), where \{$\cdots$\} denotes repetition (repeating zero or more times),
[$\cdots$] denotes option (either included or discarded),  and the terminals are doubly quoted by \mbox{\tt ``} and \mbox{\tt ''}. We stress that the syntactic rules of FIRRTL presented here is a simplification of those in the official specification.\footnote{\url{https://github.com/chipsalliance/firrtl-spec/blob/main/spec.md}}
%
For instance, indent, dedent and newline terminals are omitted for readability;
the types and constructs for analog signals are also omitted, since we focus on digital circuits in this work.
In the syntax, memory is also omitted since its semantics is still under active debate.\footnote{\url{https://github.com/chipsalliance/chisel/issues/3112}}

%\begin{itemize}
%\item A circuit in FIRRTL consists of modules, each of which include several ports and statements.
%\item Each port has an identifier and types and is either an input or output.
%\item Each statement defines a wire, register, memory, or node, connect (<=) or partial connect (<-), attach (for analog signals), or conditional (when), and so on.
%\item Expressions are widely used in statements. Each expression is an unsigned/signed integer constant, a reference, a multiplexer, a validif expression, or an expression obtained by applying primitive operations on ground types.
%\item A ground type is ``Clock'', ``Reset'', ``AsyncReset'', ``UInt'' (unsigned integer), ``SInt'' (signed integer), ``Analog'', or ``Fixed'' (fixed point number).
%\item An aggregate type is either declared by describing the (possibly nested) fields, or an array (e.g. UInt[10]).
%\end{itemize}

%%%%%%%%%%%%%%%%%%% TL this table is just for me to understand
%\begin{tabular}{c|c}
%	c & circuit\\
%	m & module \\
%	p & port\\
%	s & statement\\
%	e & expression\\
%	r & ?\\
%	w & width\\
%	t (tg, ta) & type\\
%	f & ?\\
%\end{tabular}
%%%%%%%%%%%%%%%%%%%%%%%%%%%%%%%%%%%%%%%%%%%%%%%%%%%%%%%%%%%

\begin{figure}[htbp]
{\small %\footnotesize
\[
\begin{array}{l l l}
%\mbox{\bf Circuit} &
c & ::= & \mbox{``\fnprettyll{circuit}''}\ id \mbox{``:''} \{m\}\\
%
%\mbox{\bf Module} &
m & ::= & \mbox{``\fnprettyll{module}''} \ id \{p\} \{s\}\\
%
%\mbox{\bf External Module} &
%em & ::= & \mbox{``extmodule''}\ id \{p\} [\mbox{``defname =''}\ id] \{\mbox{``parameter''}\ id \mbox{``= ''}(string \mid int)\}\\
%
%\mbox{\bf Port} &
p & ::= & (\mbox{``\fnprettyll{input}''} \mid \mbox{``\fnprettyll{output}''})\ id\mbox{``:''}t \\
%
 %\mbox{\bf Statment} &
 s & ::= & \mbox{``\fnprettyll{wire}''}\ id \mbox{``:''}t  \mid \mbox{``\fnprettyll{reg}''}\ id \mbox{``:''}t\ e [\mbox{``(\fnprettyll{with}:\{\fnprettyll{reset}=$>$(''} e\mbox{``,''} e \mbox{``)\})''}] \\
 %&
  &   &  \mid
%  \mid memory \mid
  \mbox{``\fnprettyll{inst}''}\ id\ \mbox{``\fnprettyll{of}''}\ id\ \mid \mbox{``\fnprettyll{node}''}\ id\ \mbox{``=''} id\  \mid r \mbox{``$<$=''} e  \\
  &  & \mid r \mbox{``$<$--''} e \mid r\ \mbox{``\fnprettyll{is invalid}''}  \mid \mbox{``\fnprettyll{when}''}\ e\mbox{``:''}\{s\} [\mbox{``\fnprettyll{else}:''} \{s\}] \mid \mbox{``\fnprettyll{skip}''}\\
%
%\mbox{\bf Memory} &
%memory & = &  \mbox{``mem''}, id, newline, indent, \\
 %&
% & & \mbox{``data-type''}, \mbox{``=>''}, type, newline, \\
 %&
% & & \mbox{``depth''}, \mbox{``=>''}, int, newline, \\
 %&
% & & \mbox{``read-latency''}, \mbox{``=>''}, int, newline, \\
 %&
% & & \mbox{``write-latency}, \mbox{``=>''}, int, newline,  \\
 %&
% & & \mbox{``read-under-write''}, \mbox{``=>''}, ruw, newline, \\
  %&
%  & & \{\mbox{``reader''}, \mbox{``=>''}, id, newline\}, \\
  %&
%  & & \{\mbox{``writer''}, \mbox{``=>''}, id, newline\}, \\
  %&
%  & & \{\mbox{``readwriter''}, \mbox{``=>''}, id, newline\}, dedent;\\
%
%\mbox{\bf Read Under Write Flag} &
%ruw & = & \mbox{``old''} \mid \mbox{``new''} \mid \mbox{``undefined''};\\
%
%\mbox{\bf Expression} &
  e & ::= & (\mbox{``\fnprettyll{UInt}''} | \mbox{``\fnprettyll{SInt}''}) [w] \mbox{``(''}int\mbox{``)''} \mid r \mid \mbox{``\fnprettyll{mux}(''}e\mbox{``,''} e\mbox{``,''} e\mbox{``)''} 
  %\\
   % &     &
 \mid \mbox{``\fnprettyll{validif}(''}e\mbox{``,''} e\mbox{``)''}
        %\mid \mbox{``\textcolor{blue}{read}(''}sr \mbox{``)''}
        \mid op\\
%
%&
r & ::= & sr \mid r \mbox{``[''} e \mbox{``]''}\\
%
sr & ::= & id \mid sr\mbox{``.''}id \mid sr\mbox{``[''}int\mbox{``]''}\\
%&
w & ::= & \mbox{``$<$''} int \mbox{``$>$''}\\
%
%&
%binarypoint & = & \mbox{``<''}, \mbox{``<''}, int, \mbox{``>''}, \mbox{``>''};\\
%
 %\mbox{\bf Bundle Field} &
%
%\mbox{\bf Data Types} &
t & ::= & tg \mid ta \\
 %&
 tg & ::= &  \mbox{``\fnprettyll{Clock}''} \mid \mbox{``\fnprettyll{Reset}''} \mid \mbox{``\fnprettyll{AsyncReset}''} \mid (\mbox{``\fnprettyll{UInt}''} \mid \mbox{``\fnprettyll{SInt}''}) [w] \\
 %&
% & & \mid \mbox{``Fixed''}, [width], [binarypoint];\\
 %&
%
  ta & ::= &\mbox{``\{''} f \{f\}\mbox{``\}''} \mid t \mbox{``[''} int \mbox{``]''} \\
%
    f & ::= & [\mbox{``\fnprettyll{flip}''}]\ id\mbox{``:''}t\\
%& & & \mbox{``UInt''}, [\mbox{``<''}, int, \mbox{``>''}] \mid \mbox{``SInt''}, [\mbox{``<''}, int, \mbox{``>''} ] \\
%& & & \mid \mbox{``Fixed''}, [\mbox{``<''}, int, \mbox{``>''} ] [\mbox{``<}\mbox{<''}, int, \mbox{``>}\mbox{>''}]\\
% &  &  & \mid \mbox{``Clock''} \mid \mbox{``Analog''},  [\mbox{``<''}, int, \mbox{``>''}] \mid \mbox{``\{''}, field, \mbox{``\}''} \mid type, \mbox{``[''}, int, \mbox{``]''}; \\
%
id & ::= & ( \mbox{``\_''} \mid l) \{ \mbox{``\_''} \mid l \mid d\} \\
%\mbox{\bf Primitive Operation} &
op & ::= & op\_2e \mid op\_1e \mid op\_1e1i \mid op\_1e2i \\
%&
op\_2e & ::= & op\_2e\_key \mbox{``(''}  e \mbox{``,''}  e \mbox{``)''} \\
%
%&
op\_2e\_key & ::= &  \mbox{``\fnprettyll{add}''}  \mid  \mbox{``\fnprettyll{sub}''}  \mid \mbox{``\fnprettyll{mul}''} \mid \mbox{``\fnprettyll{div}''} \mid \mbox{``\fnprettyll{mod}''} \mid \mbox{``\fnprettyll{lt}''} \mid \mbox{``\fnprettyll{leq}''} \mid \mbox{``\fnprettyll{gt}''} \mid \mbox{``\fnprettyll{geq}''}\\
%&
& &   \mid \mbox{``\fnprettyll{eq}''} \mid \mbox{``\fnprettyll{neq}''}  \mid \mbox{``\fnprettyll{dshl}''}  \mid \mbox{``\fnprettyll{dshr}''} \mid \mbox{``\fnprettyll{and}''} \mid \mbox{``\fnprettyll{or}''} \mid \mbox{``\fnprettyll{xor}''} \mid \mbox{``\fnprettyll{cat}''} \\
 %
%&
op\_1e & ::= & op\_1e\_key \mbox{``(''}  e \mbox{``)''} \\
%&
op\_1e\_key & ::= & \mbox{``\fnprettyll{asUInt}''} \mid \mbox{``\fnprettyll{asSInt}''} \mid \mbox{``\fnprettyll{asClock}''} \mid \mbox{``\fnprettyll{cvt}''} \mid \mbox{``\fnprettyll{neg}''} \mid \mbox{``\fnprettyll{not}''}  \mid \mbox{``\fnprettyll{andr}''} \mid \mbox{``\fnprettyll{orr}''} \mid  \mbox{``\fnprettyll{xorr}''}\\
%
%&
op\_1e1i & ::= & op\_1e1i\_key \mbox{``(''}  e \mbox{``,''} int \mbox{``)''} \\
%&
op\_1e1i\_key & ::= & \mbox{``\fnprettyll{pad}''} \mid \mbox{``\fnprettyll{shl}''} \mid \mbox{``\fnprettyll{shr}''} \mid \mbox{``\fnprettyll{head}''} \mid \mbox{``\fnprettyll{tail}''} \\
%&
op\_1e2i & ::= & op\_1e2i\_key \mbox{``(''}  e \mbox{``,''}  int  \mbox{``,''}  int \mbox{``)''} \\
%&
op\_1e2i\_key & ::= & \mbox{``\fnprettyll{bits}''}\\
%
int & ::= & [\mbox{``-''}] d \{d\} \mid \mbox{``{\tt "}b''} [\mbox{``-''}] d\_b \{d\_b\} \mbox{``{\tt "}''} \\
& & \mid \mbox{``{\tt "}o''} [\mbox{``-''}] d\_o \{d\_o\} \mbox{``{\tt "}''} \mid
\mbox{``{\tt "}h''} [\mbox{``-''}] d\_h \{d\_h\} \mbox{``{\tt "}''} \\
%
%d & ::= &\mbox{``0''} \mid  \mbox{ ``1''}\mid \mbox{``2''} \mid\mbox{``3''} \mid  \mbox{ ``4''}\mid \mbox{``5''} \mid\mbox{``6''} \mid  \mbox{ ``7''}\mid \mbox{``8''} \mid\mbox{``9''} \\
%d\_b & ::= &\mbox{``0''} \mid  \mbox{ ``1''}  \\
%d\_o & ::= &\mbox{``0''} \mid  \mbox{ ``1''}\mid \mbox{``2''} \mid\mbox{``3''} \mid  \mbox{ ``4''}\mid \mbox{``5''} \mid\mbox{``6''} \mid  \mbox{ ``7''} \\
%d\_h & ::= &\mbox{``0''} \mid  \mbox{ ``1''}\mid \mbox{``2''} \mid\mbox{``3''} \mid  \mbox{ ``4''}\mid \mbox{``5''} \mid\mbox{``6''} \mid  \mbox{ ``7''}\mid \mbox{``8''} \mid\mbox{``9''}  \mid  \mbox{ ``A''}\mid \mbox{``B''} \mid\mbox{``C''} \mid  \mbox{ ``D''}\mid \mbox{``E''} \mid\mbox{``F''}\\
%

%\mbox{\bf File Information Token} &
%info & = & \mbox{``@''}, \mbox{``[''}, \{string, \mbox{`` ''}, linecol\}, \mbox{``]''};\\
%&
%linecol & = &  digit\_dec, \{digit\_dec\}, \mbox{``:''}, digit\_dec, \{digit\_dec\};\\
%&
%
\end{array}
\]
$d$ (resp. $d\_b$, $d\_o$ and $d\_h$) denotes a decimal (resp. binary, octal and hexadecimal) digit, and $l$ denotes a letter
}
\caption{Syntax of FIRRTL}
\label{fig-firrtl-synt-simp}%\vspace{-2mm}
\end{figure}

\hide{
\begin{figure}[htbp]
{\small %\footnotesize
\[
\begin{array}{l l l}
%\mbox{\bf Circuit} &
c & ::= & \mbox{\tt ``\fnprettyll{circuit}''}\ id\mbox{\tt ``:''} \{m\}\\
%
%\mbox{\bf Module} &
m & ::= & \mbox{\tt ``\fnprettyll{module}''} \ id \{p\} \{s\}\\
%
%\mbox{\bf External Module} &
%em & ::= & \mbox{\tt ``extmodule''}\ id \{p\} [\mbox{\tt ``defname =''}\ id] \{\mbox{\tt ``parameter''}\ id \mbox{\tt ``= ''}(string \mid int)\}\\
%
%\mbox{\bf Port} &
p & ::= & (\mbox{\tt ``\fnprettyll{input}''} \mid \mbox{\tt ``\fnprettyll{output}''})\ id\mbox{\tt ``:''}t \\
%
 %\mbox{\bf Statment} &
 s & ::= & \mbox{\tt ``\fnprettyll{wire}''}\ id\mbox{\tt ``:''}t  \mid \mbox{\tt ``\fnprettyll{reg}''}\ id\mbox{\tt ``:''}t\ e [\mbox{\tt ``(\fnprettyll{with}:\{\fnprettyll{reset}=>(''} e\mbox{\tt ``,''} e \mbox{\tt ``)\})''}] \\
 %&
  &   &  \mid
%  \mid memory \mid
  \mbox{\tt ``\fnprettyll{inst}''}\ id\ \mbox{\tt ``\fnprettyll{of}''}\ id \mid \mbox{\tt ``\fnprettyll{node}''}\ id \mbox{\tt ``=''} id  \mid r \mbox{\tt ``<=''} e  \\
  &  & \mid r \mbox{\tt ``<-''} e \mid r\ \mbox{\tt ``\fnprettyll{is invalid}''}  \mid \mbox{\tt ``\fnprettyll{when}''}\ e\mbox{\tt ``:''}\{s\} [\mbox{\tt ``\fnprettyll{else}:''} \{s\}] \mid \mbox{\tt ``\fnprettyll{skip}''}\\
%
%\mbox{\bf Memory} &
%memory & = &  \mbox{``mem''}, id, newline, indent, \\
 %&
% & & \mbox{``data-type''}, \mbox{``=>''}, type, newline, \\
 %&
% & & \mbox{``depth''}, \mbox{``=>''}, int, newline, \\
 %&
% & & \mbox{``read-latency''}, \mbox{``=>''}, int, newline, \\
 %&
% & & \mbox{``write-latency}, \mbox{``=>''}, int, newline,  \\
 %&
% & & \mbox{``read-under-write''}, \mbox{``=>''}, ruw, newline, \\
  %&
%  & & \{\mbox{``reader''}, \mbox{``=>''}, id, newline\}, \\
  %&
%  & & \{\mbox{``writer''}, \mbox{``=>''}, id, newline\}, \\
  %&
%  & & \{\mbox{``readwriter''}, \mbox{``=>''}, id, newline\}, dedent;\\
%
%\mbox{\bf Read Under Write Flag} &
%ruw & = & \mbox{``old''} \mid \mbox{``new''} \mid \mbox{``undefined''};\\
%
%\mbox{\bf Expression} &
  e & ::= & (\mbox{\tt ``\fnprettyll{UInt}''} | \mbox{\tt ``\fnprettyll{SInt}''}) [w] \mbox{\tt ``(''}int\mbox{\tt ``)''} \mid r \mid \mbox{\tt ``\fnprettyll{mux}(''}e\mbox{\tt ``,''} e\mbox{\tt ``,''} e\mbox{\tt ``)''} \\
  %\\
    &     &
 \mid \mbox{\tt ``\fnprettyll{validif}(''}e\mbox{\tt ``,''} e\mbox{\tt ``)''}
        %\mid \mbox{\tt ``\textcolor{blue}{read}(''}sr \mbox{\tt ``)''}
        \mid op\\
%
%&
r & ::= & sr \mid r \mbox{\tt ``[''} e \mbox{\tt ``]''}\\
%
sr & ::= & id \mid sr\mbox{\tt ``.''}id \mid sr\mbox{\tt ``[''}int\mbox{\tt ``]''}\\
%&
w & ::= & \mbox{\tt ``<''} int \mbox{\tt ``>''}\\
%
%&
%binarypoint & = & \mbox{``<''}, \mbox{``<''}, int, \mbox{``>''}, \mbox{``>''};\\
%
 %\mbox{\bf Bundle Field} &
%
%\mbox{\bf Data Types} &
t & ::= & tg \mid ta \\
 %&
 tg & ::= &  \mbox{\tt ``\fnprettyll{Clock}''} \mid \mbox{\tt ``\fnprettyll{Reset}''} \mid \mbox{\tt ``\fnprettyll{AsyncReset}''} \mid (\mbox{\tt ``\fnprettyll{UInt}''} \mid \mbox{\tt ``\fnprettyll{SInt}''}) [w] \\
 %&
% & & \mid \mbox{``Fixed''}, [width], [binarypoint];\\
 %&
%
  ta & ::= &\mbox{\tt ``\{''} f \{f\}\mbox{\tt ``\}''} \mid t \mbox{\tt ``[''} int \mbox{\tt ``]''} \\
%
    f & ::= & [\mbox{\tt ``\fnprettyll{flip}''}]\ id\mbox{\tt ``:''}t\\
%& & & \mbox{``UInt''}, [\mbox{``<''}, int, \mbox{``>''}] \mid \mbox{``SInt''}, [\mbox{``<''}, int, \mbox{``>''} ] \\
%& & & \mid \mbox{``Fixed''}, [\mbox{``<''}, int, \mbox{``>''} ] [\mbox{``<}\mbox{<''}, int, \mbox{``>}\mbox{>''}]\\
% &  &  & \mid \mbox{``Clock''} \mid \mbox{``Analog''},  [\mbox{``<''}, int, \mbox{``>''}] \mid \mbox{``\{''}, field, \mbox{``\}''} \mid type, \mbox{``[''}, int, \mbox{``]''}; \\
%
id & ::= & ( \mbox{\tt ``\_''} \mid l) \{ \mbox{\tt ``\_''} \mid l \mid d\} \\
%\mbox{\bf Primitive Operation} &
op & ::= & op\_2e \mid op\_1e \mid op\_1e1i \mid op\_1e2i \\
%&
op\_2e & ::= & op\_2e\_key \mbox{\tt ``(''}  e \mbox{\tt ``,''}  e \mbox{\tt ``)''} \\
%
%&
op\_2e\_key & ::= &  \mbox{\tt ``\fnprettyll{add}''}  \mid  \mbox{\tt ``\fnprettyll{sub}''}  \mid \mbox{\tt ``\fnprettyll{mul}''} \mid \mbox{\tt ``\fnprettyll{div}''} \mid \mbox{\tt ``\fnprettyll{mod}''} \mid \mbox{\tt ``\fnprettyll{lt}''} \mid \mbox{\tt ``\fnprettyll{leq}''} \mid \mbox{\tt ``\fnprettyll{gt}''} \mid \mbox{\tt ``\fnprettyll{geq}''}\\
%&
& &   \mid \mbox{\tt ``\fnprettyll{eq}''} \mid \mbox{\tt ``\fnprettyll{neq}''}  \mid \mbox{\tt ``\fnprettyll{dshl}''}  \mid \mbox{\tt ``\fnprettyll{dshr}''} \mid \mbox{\tt ``\fnprettyll{and}''} \mid \mbox{\tt ``\fnprettyll{or}''} \mid \mbox{\tt ``\fnprettyll{xor}''} \mid \mbox{\tt ``\fnprettyll{cat}''} \\
 %
%&
op\_1e & ::= & op\_1e\_key \mbox{\tt ``(''}  e \mbox{\tt ``)''} \\
%&
op\_1e\_key & ::= & \mbox{\tt ``\fnprettyll{asUInt}''} \mid \mbox{\tt ``\fnprettyll{asSInt}''} \mid \mbox{\tt ``\fnprettyll{asClock}''} \mid \mbox{\tt ``\fnprettyll{cvt}''} \mid \mbox{\tt ``\fnprettyll{neg}''} \mid \mbox{\tt ``\fnprettyll{not}''}  \mid \mbox{\tt ``\fnprettyll{andr}''} \mid \mbox{\tt ``\fnprettyll{orr}''} \mid  \mbox{\tt ``\fnprettyll{xorr}''}\\
%
%&
op\_1e1i & ::= & op\_1e1i\_key \mbox{\tt ``(''}  e \mbox{\tt ``,''} int \mbox{\tt ``)''} \\
%&
op\_1e1i\_key & ::= & \mbox{\tt ``\fnprettyll{pad}''} \mid \mbox{\tt ``\fnprettyll{shl}''} \mid \mbox{\tt ``\fnprettyll{shr}''} \mid \mbox{\tt ``\fnprettyll{head}''} \mid \mbox{\tt ``\fnprettyll{tail}''} \\
%&
op\_1e2i & ::= & op\_1e2i\_key \mbox{\tt ``(''}  e \mbox{\tt ``,''}  int  \mbox{\tt ``,''}  int \mbox{\tt ``)''} \\
%&
op\_1e2i\_key & ::= & \mbox{\tt ``\fnprettyll{bits}''}\\
%
int & ::= & [\mbox{\tt ``-''}] d \{d\} \mid \mbox{\tt ``"b''} [\mbox{\tt ``-''}] d\_b \{d\_b\} \mbox{\tt ``"''} \\
& & \mid \mbox{\tt ``"o''} [\mbox{\tt ``-''}] d\_o \{d\_o\} \mbox{\tt ``"''} \mid
\mbox{\tt ``"h''} [\mbox{\tt ``-''}] d\_h \{d\_h\} \mbox{\tt ``"''} \\
%
%d & ::= &\mbox{\tt ``0''} \mid  \mbox{\tt  ``1''}\mid \mbox{\tt ``2''} \mid\mbox{\tt ``3''} \mid  \mbox{\tt  ``4''}\mid \mbox{\tt ``5''} \mid\mbox{\tt ``6''} \mid  \mbox{\tt  ``7''}\mid \mbox{\tt ``8''} \mid\mbox{\tt ``9''} \\
%d\_b & ::= &\mbox{\tt ``0''} \mid  \mbox{\tt  ``1''}  \\
%d\_o & ::= &\mbox{\tt ``0''} \mid  \mbox{\tt  ``1''}\mid \mbox{\tt ``2''} \mid\mbox{\tt ``3''} \mid  \mbox{\tt  ``4''}\mid \mbox{\tt ``5''} \mid\mbox{\tt ``6''} \mid  \mbox{\tt  ``7''} \\
%d\_h & ::= &\mbox{\tt ``0''} \mid  \mbox{\tt  ``1''}\mid \mbox{\tt ``2''} \mid\mbox{\tt ``3''} \mid  \mbox{\tt  ``4''}\mid \mbox{\tt ``5''} \mid\mbox{\tt ``6''} \mid  \mbox{\tt  ``7''}\mid \mbox{\tt ``8''} \mid\mbox{\tt ``9''}  \mid  \mbox{\tt  ``A''}\mid \mbox{\tt ``B''} \mid\mbox{\tt ``C''} \mid  \mbox{\tt  ``D''}\mid \mbox{\tt ``E''} \mid\mbox{\tt ``F''}\\
%

%\mbox{\bf File Information Token} &
%info & = & \mbox{``@''}, \mbox{``[''}, \{string, \mbox{`` ''}, linecol\}, \mbox{``]''};\\
%&
%linecol & = &  digit\_dec, \{digit\_dec\}, \mbox{``:''}, digit\_dec, \{digit\_dec\};\\
%&
%
\end{array}
\]
$d$ (resp. $d\_b$, $d\_o$ and $d\_h$) denotes a decimal (resp. binary, octal and hexadecimal) digit, and $l$ denotes a letter
}
\caption{Syntax of FIRRTL}
\label{fig-firrtl-synt-simp}%\vspace{-2mm}
\end{figure}
}

\begin{comment}
\begin{figure}[htbp]
{\footnotesize
\[
\begin{array}{l l l}
%\mbox{\bf Circuit} &
c & ::= & \mbox{\tt ``circuit''}\ id\mbox{\tt ``:''} \{m\}\\
%
%\mbox{\bf Module} &
m & ::= & \mbox{\tt ``module''} \ id \{p\} \{s\}\\
%
%\mbox{\bf External Module} &
%em & ::= & \mbox{\tt ``extmodule''}\ id \{p\} [\mbox{\tt ``defname =''}\ id] \{\mbox{\tt ``parameter''}\ id \mbox{\tt ``= ''}(string \mid int)\}\\
%
%\mbox{\bf Port} &
p & ::= & (\mbox{\tt ``input''} \mid \mbox{\tt ``output''})\ id\mbox{\tt ``:''}t \\
%
 %\mbox{\bf Statment} &
 s & ::= & \mbox{\tt ``wire''}\ id\mbox{\tt ``:''}t  \mid \mbox{\tt ``reg''}\ id\mbox{\tt ``:''}t\ e [\mbox{\tt ``(with:\{reset=>(''} e\mbox{\tt ``,''} e \mbox{\tt ``)\})''}] \\
 %&
  &   &  \mid
%  \mid memory \mid
  \mbox{\tt ``inst''}\ id\ \mbox{\tt ``of''}\ id \mid \mbox{\tt ``node''}\ id \mbox{\tt ``=''} id  \mid r \mbox{\tt ``<=''} e  \mid r \mbox{\tt ``<-''} e \mid r\ \mbox{\tt ``is invalid''} \\
%&
&   &\mid \mbox{\tt ``when''}\ e\mbox{\tt ``:''}\{s\} [\mbox{\tt ``else:''} \{s\}] \mid \mbox{\tt ``skip''}\\
%
%\mbox{\bf Memory} &
%memory & = &  \mbox{``mem''}, id, newline, indent, \\
 %&
% & & \mbox{``data-type''}, \mbox{``=>''}, type, newline, \\
 %&
% & & \mbox{``depth''}, \mbox{``=>''}, int, newline, \\
 %&
% & & \mbox{``read-latency''}, \mbox{``=>''}, int, newline, \\
 %&
% & & \mbox{``write-latency}, \mbox{``=>''}, int, newline,  \\
 %&
% & & \mbox{``read-under-write''}, \mbox{``=>''}, ruw, newline, \\
  %&
%  & & \{\mbox{``reader''}, \mbox{``=>''}, id, newline\}, \\
  %&
%  & & \{\mbox{``writer''}, \mbox{``=>''}, id, newline\}, \\
  %&
%  & & \{\mbox{``readwriter''}, \mbox{``=>''}, id, newline\}, dedent;\\
%
%\mbox{\bf Read Under Write Flag} &
%ruw & = & \mbox{``old''} \mid \mbox{``new''} \mid \mbox{``undefined''};\\
%
%\mbox{\bf Expression} &
e & ::= & (\mbox{\tt ``UInt''} | \mbox{\tt ``SInt''}) [w] \mbox{\tt ``(''}int\mbox{\tt ``)''} \mid r \mid \mbox{\tt ``mux(''}e\mbox{\tt ``,''} e\mbox{\tt ``,''} e\mbox{\tt ``)''}
%  \mid \mbox{\tt ``\prettyll{validif}(''}e\mbox{\tt ``,''} e\mbox{\tt ``)''} \mid \mbox{\tt ``read(''}sr \mbox{\tt ``)''}
  \mid op\\
%
%&
r & ::= & sr \mid r \mbox{\tt ``[''} e \mbox{\tt ``]''}\\
%
sr & ::= & id \mid sr\mbox{\tt ``.''}id \mid sr\mbox{\tt ``[''}int\mbox{\tt ``]''}\\
%&
w & ::= & \mbox{\tt ``<''} int \mbox{\tt ``>''}\\
%
%&
%binarypoint & = & \mbox{``<''}, \mbox{``<''}, int, \mbox{``>''}, \mbox{``>''};\\
%
 %\mbox{\bf Bundle Field} &
%
%\mbox{\bf Data Types} &
t & ::= & tg \mid ta \\
 %&
 tg & ::= &  \mbox{\tt ``Clock''} \mid \mbox{\tt ``Reset''} \mid \mbox{\tt ``AsyncReset''} \mid (\mbox{\tt ``UInt''} \mid \mbox{\tt ``SInt''}) [w] \\
 %&
% & & \mid \mbox{``Fixed''}, [width], [binarypoint];\\
 %&
%
  ta & ::= &\mbox{\tt ``\{''} f \{f\}\mbox{\tt ``\}''} \mid t \mbox{``[''} int \mbox{``]''} \\
%
    f & ::= & [\mbox{\tt ``flip''}]\ id\mbox{\tt ``:''}t\\
%& & & \mbox{``UInt''}, [\mbox{``<''}, int, \mbox{``>''}] \mid \mbox{``SInt''}, [\mbox{``<''}, int, \mbox{``>''} ] \\
%& & & \mid \mbox{``Fixed''}, [\mbox{``<''}, int, \mbox{``>''} ] [\mbox{``<}\mbox{<''}, int, \mbox{``>}\mbox{>''}]\\
% &  &  & \mid \mbox{``Clock''} \mid \mbox{``Analog''},  [\mbox{``<''}, int, \mbox{``>''}] \mid \mbox{``\{''}, field, \mbox{``\}''} \mid type, \mbox{``[''}, int, \mbox{``]''}; \\
%
id & ::= & ( \mbox{\tt ``\_''} \mid l) \{ \mbox{\tt ``\_''} \mid l \mid d\} \\
%\mbox{\bf Primitive Operation} &
op & ::= & op\_2e \mid op\_1e \mid op\_1e1i \mid op\_1e2i \\
%&
op\_2e & ::= & op\_2e\_key \mbox{\tt ``(''}  e \mbox{\tt ``,''}  e \mbox{\tt ``)''} \\
%
%&
op\_2e\_key & ::= &  \mbox{\tt ``add''}  \mid  \mbox{\tt ``sub''}  \mid \mbox{\tt ``mul''} \mid \mbox{\tt ``div''} \mid \mbox{\tt ``mod''} \mid \mbox{\tt ``lt''} \mid \mbox{\tt ``leq''} \mid \mbox{\tt ``gt''} \mid \mbox{\tt ``geq''} \mid \mbox{\tt ``eq''} \mid \mbox{\tt ``neq''} \\
%&
& &   \mid \mbox{\tt ``dshl''}  \mid \mbox{\tt ``dshr''} \mid \mbox{\tt ``and''} \mid \mbox{\tt ``or''} \mid \mbox{\tt ``xor''} \mid \mbox{\tt ``cat''} \\
 %
%&
op\_1e & ::= & op\_1e\_key \mbox{\tt ``(''}  e \mbox{\tt ``)''} \\
%&
op\_1e\_key & ::= & \mbox{\tt ``asUInt''} \mid \mbox{\tt ``asSInt''} \mid \mbox{\tt ``asClock''} \mid \mbox{\tt ``cvt''} \mid \mbox{\tt ``neg''} \mid \mbox{\tt ``not''}  \mid \mbox{\tt ``andr''} \mid \mbox{\tt ``orr''} \mid  \mbox{\tt ``xorr''}\\
%
%&
op\_1e1i & ::= & op\_1e1i\_key \mbox{\tt ``(''}  e \mbox{\tt ``,''} int \mbox{\tt ``)''} \\
%&
op\_1e1i\_key & ::= & \mbox{\tt ``pad''} \mid \mbox{\tt ``shl''} \mid \mbox{\tt ``shr''} \mid \mbox{\tt ``head''} \mid \mbox{\tt ``tail''} \\
%&
op\_1e2i & ::= & op\_1e2i\_key \mbox{\tt ``(''}  e \mbox{\tt ``,''}  int  \mbox{\tt ``,''}  int \mbox{\tt ``)''} \\
%&
op\_1e2i\_key & ::= & \mbox{\tt ``bits''}\\
%
int & ::= & [\mbox{\tt ``-''}] d \{d\} \mid \mbox{\tt ``"b''} [\mbox{\tt ``-''}] d\_b \{d\_b\} \mbox{\tt ``"''}
 \mid \mbox{\tt ``"o''} [\mbox{\tt ``-''}] d\_o \{d\_o\} \mbox{\tt ``"''} \mid
\mbox{\tt ``"h''} [\mbox{\tt ``-''}] d\_h \{d\_h\} \mbox{\tt ``"''}
%
%\mbox{\bf File Information Token} &
%info & = & \mbox{``@''}, \mbox{``[''}, \{string, \mbox{`` ''}, linecol\}, \mbox{``]''};\\
%&
%linecol & = &  digit\_dec, \{digit\_dec\}, \mbox{``:''}, digit\_dec, \{digit\_dec\};\\
%&
%
\end{array}
\]
}
\caption{Syntax of FIRRTL}
\label{fig-firrtl-synt-simp}
\end{figure}
\end{comment}

According to the syntax, FIRRTL features high-level constructs such as vector types, bundle types, conditional statements, connects, partial connects, and modules.
These high-level constructs are then gradually removed by a sequence of \emph{lowering transformations}. During each lowering transformation, the circuit is rewritten into an equivalent circuit using simpler, lower-level
constructs. Eventually the circuit is simplified to its most restricted form, resembling a structured netlist,
which allows for easy translation to an output language (e.g. Verilog).
This form is given the name lowered FIRRTL (abbreviated as LoFIRRTL) and is a strict subset of the full FIRRTL language.
For clarity, the original form is referred to as high FIRRTL (HiFIRRTL).
\footnote{Besides LoFIRRTL, there is another lowered form of FIRRTL, referred to as MidFIRRTL, which resides between HiFIRRTL and LoFIRRTL. We do not introduce MidFIRRTL in this work,
since it can be seen as an internal representation of a FIRRTL compiler that does not interface with the other parts of the Chisel toolchain.}


A FIRRTL circuit is a LoFIRRTL circuit if it obeys the following restrictions:
\begin{itemize}
\item All widths are explicitly defined.
\item No conditional statements (\prettyll{when}) are used.
\item All circuit components are declared with a ground type.
\item No partial connect statements (\prettyll{<-}) are used.
\item Each component is connected to at most once.
\end{itemize}
More precisely, a LoFIRRTL circuit is a FIRRTL circuit that is defined by the syntactic rules in Figure~\ref{fig-lofirrtl-synt}
and satisfies that in each module, each component is connected to at most once.
\begin{figure}[htbp]
{\small
\[
\begin{array}{l l l}
 s & ::= & \mbox{\tt ``\fnprettyll{wire}''}\ id\mbox{\tt ``:''}t  \mid \mbox{\tt ``\fnprettyll{reg}''}\ id\mbox{\tt ``:''}t\ e [\mbox{\tt ``(\fnprettyll{with}:\{\fnprettyll{reset}=>(''} e\mbox{\tt ``,''} e \mbox{\tt ``)\})''}] \\
 %&
  &   &  \mid
%  \mid memory \mid
  \mbox{\tt ``\fnprettyll{inst}''}\ id\ \mbox{\tt ``\fnprettyll{of}''}\ id \mid \mbox{\tt ``\fnprettyll{node}''}\ id \mbox{\tt ``=''} id  \mid r \mbox{\tt ``<=''} e  \mid r\ \mbox{\tt ``\fnprettyll{is invalid}''}\mid \mbox{\tt ``\fnprettyll{skip}''}\\
%
%\mbox{\bf Expression} &
e & ::= & (\mbox{\tt ``\fnprettyll{UInt}''} | \mbox{\tt ``\fnprettyll{SInt}''})\ w \mbox{\tt ``(''}int\mbox{\tt ``)''} \mid r \mid \mbox{\tt ``\fnprettyll{mux}(''}e\mbox{\tt ``,''} e\mbox{\tt ``,''} e\mbox{\tt ``)''}
% &  & \mid \mbox{\tt ``\fnprettyll{validif}(''}e\mbox{\tt ``,''} e\mbox{\tt ``)''} \mid \mbox{\tt ``\textcolor{blue}{read}(''}sr \mbox{\tt ``)''}
 \mid op\\
%
%&
r & ::= & sr \mid r \mbox{\tt ``[''} e \mbox{\tt ``]''}\\
%
sr & ::= & id \mid sr\mbox{\tt ``.''}id \mid sr\mbox{\tt ``[''}int\mbox{\tt ``]''}\\
%&
w & ::= & \mbox{\tt ``<''} int \mbox{\tt ``>''}\\
%
%&
%binarypoint & = & \mbox{``<''}, \mbox{``<''}, int, \mbox{``>''}, \mbox{``>''};\\
%
 %\mbox{\bf Bundle Field} &
%
%\mbox{\bf Data Types} &
t & ::= & tg \\
 %&
 tg & ::= &  \mbox{\tt ``\fnprettyll{Clock}''} \mid \mbox{\tt ``\fnprettyll{AsyncReset}''} \mid (\mbox{\tt ``\fnprettyll{UInt}''} \mid \mbox{\tt ``\fnprettyll{SInt}''})\ w
 %&
% & & \mid \mbox{``Fixed''}, [width], [binarypoint];\\
 %&
%
%\mbox{\bf File Information Token} &
%info & = & \mbox{``@''}, \mbox{``[''}, \{string, \mbox{`` ''}, linecol\}, \mbox{``]''};\\
%&
%linecol & = &  digit\_dec, \{digit\_dec\}, \mbox{``:''}, digit\_dec, \{digit\_dec\};\\
%&
%
\end{array}
\]
}
\caption{Syntax of LoFIRRTL (Rules that are identical to those in Figure~\ref{fig-firrtl-synt-simp} are not repeated)}
\label{fig-lofirrtl-synt}
\end{figure}

Note that in Figure~\ref{fig-lofirrtl-synt}, the rules for $f$, $id$, $op$ are the same as those in Figure~\ref{fig-firrtl-synt-simp}  (and thus omitted). %to avoid tediousness.
In a nutshell, the syntactic rules in Figure~\ref{fig-lofirrtl-synt} are obtained from those in Figure~\ref{fig-firrtl-synt-simp} by replacing $[w]$ with $w$ (so that the width becomes mandatory),
removing the \prettyll{when} statement and the partial connect statement (\prettyll{<-}) from the rule for $s$,
and removing the rule for aggregate types $ta$.
%
The requirement that each component is connected to at most once
gives LoFIRRTL a flavor similar to static single-assignment (SSA) form of software programs.
Programs in SSA form are those where each variable is assigned at most once and each variable is defined before it is used.
For instance, the program in Figure~\ref{fig-lofirrtl-exmp}(a) is a LoFIRRTL program,
while the program in Figure~\ref{fig-lofirrtl-exmp}(b) is not,
since the output port \prettyll{myport1} is connected to twice from the input ports \prettyll{a} and \prettyll{b}, respectively.
\begin{figure}[htbp]
\centering
\begin{subfigure}[b]{0.4\textwidth}
         \centering
%[basicstyle=\footnotesize\ttfamily]
\begin{lstlisting}
module MyModule :
  input a: UInt<1>
  input b: UInt<1>
  output myport1: UInt<1>
  output myport2: UInt<1>
  myport1 <= b
  myport2 <= a
\end{lstlisting}
\vspace*{0.92\baselineskip}
\caption{}
\end{subfigure}
\hspace{4mm}
%     \hfill
\begin{subfigure}[b]{0.4\textwidth}
         \centering
%[basicstyle=\footnotesize\ttfamily]
\begin{lstlisting}
module MyModule :
  input a: UInt<1>
  input b: UInt<1>
  output myport1: UInt<1>
  output myport2: UInt<1>
  myport1 <= a
  myport1 <= b
  myport2 <= a
\end{lstlisting}
\caption{}
\end{subfigure}
\caption{Examples for illustrating the syntactic restrictions of LoFIRRTL}
\label{fig-lofirrtl-exmp}
\end{figure}

%%%%%%%%%%%%%%%%%%%%%%%%%%%%%%%%%%%%%%%%%%%%%%%%%%%%%%%%%%%%%%%%%%%

\subsection{Lowering Transformations of FIRRTL}
\label{subsec-lowering-transformations}

In the sequel, we give an informal introduction to the lowering transformations (aka compilation passes) of  FIRRTL.

%The dependencies between these compilation passes are illustrated in Figure~\ref{fig-c-pass} \tl{redo this}.


%\begin{figure}[htbp]
%\centering
%\includegraphics[width = 0.65\textwidth]{firrtl-compilation-passes}
%\caption{Dependencies between FIRRTL compilation passes}
%\label{fig-c-pass}
%\end{figure}


%% \paragraph*{ResolveKinds} xxx \zhilin{@Xiaomu, please write an introduction here.}
%
%\paragraph*{InferTypes} The InferTypes pass infers the types of components. For instance, InferTypes infers from the code snippet
%\prettyll{ wire x: UInt   x <= UInt<5>(15) x <= UInt<1>(1) } that the type of the wire $x$ is \prettyll{ UInt}. Note that for a connect statement   \prettyll{ x <= e}, InferTypes only checks that the type of
% the expression $e$ is equivalent to that of the component $x$ (where the bit widths are ignored), and does not update the bit width of $x$ to that of the expression $e$.

\tl{incorporate infertypes please}

\paragraph*{InferWidths}
The InferWidths pass infers the widths of all components whose widths are unknown. %, after the InferTypes pass has been applied \tl{to be revised}.
It is a \emph{global} width inference algorithm. Instances of the same module with unknown input port widths will be assigned the largest width of all assignments to each of its instance ports. It infers the smallest width that is larger than all assigned widths to a component. Note that this means that dummy assignments that are overwritten by last-connect semantics can still influence width inference. The last-connect semantics means that the last connect statement for a component will take effects and all the preceding connect statements for it will be ignored  in the end. For instance, InferWidths infers from the code snippet \prettyll{wire x: UInt   x <= UInt<5>(15)    x <= UInt<1>(1)} that the width of the wire $x$ is $5$, but with a value \prettyll{UInt<1>(1)}.


%\paragraph*{InferResets} %% that infers concrete resets, i.e. AsyncReset or UInt<1>,
%The InferResets pass infers the concrete reset types for the abstract reset type \prettyll{Reset}, namely, asynchronous resets represented by  \prettyll{AsyncReset} or synchronous resets represented by \prettyll{UInt<1>}. For example, in the code snippet \prettyll{ input in: UInt<1>   wire rst: Reset   rst <= in } with an input port \prettyll{in} of type \prettyll{UInt<1>} and a wire \prettyll{rst} of abstract type \prettyll{Reset}, the \prettyll{rst} wire will be inferred as a synchronous reset based on the type of the input port \prettyll{in}. If a component of \prettyll{Reset} type is instantiated with different concrete reset types, then a conflict occurs and an error is reported.

\paragraph{ExpandConnects} The ExpandConnects pass expands the connect and partial connect statements, after the types of components have been inferred by the InferTypes pass.
For instance, the ExpandConnects pass transforms the program shown in Figure~\ref{fig-lofirrtl-exmp2}(a) into the one shown in Figure~\ref{fig-lofirrtl-exmp2}(b).


\begin{figure}[h]
\centering
\begin{subfigure}[b]{0.4\textwidth}
         \centering
%[basicstyle=\footnotesize\ttfamily]
\begin{lstlisting}[boxpos=t,aboveskip=0pt]
module MyModule:
  input x: {a: UInt, b: UInt}
  output y:{a: UInt, b: UInt}
  y <= x
\end{lstlisting}
\vspace*{0.92\baselineskip}
\caption{}
\end{subfigure}
\hspace{4mm}
%     \hfill
\begin{subfigure}[b]{0.4\textwidth}
         \centering
%[basicstyle=\footnotesize\ttfamily]
\begin{lstlisting}
 module MyModule:
  input x: {a: UInt, b: UInt}
  output y:{a: UInt, b: UInt}
  y.a <= x.a
  y.b <= x.b
\end{lstlisting}
\caption{}
\end{subfigure}
\caption{Examples for illustrating the ExpandConnects pass}
\label{fig-lofirrtl-exmp2}
\end{figure}

%
%\begin{figure}[h]
%\centering
%\begin{tabular}{cccc}
%%[basicstyle=\footnotesize\ttfamily]
%\begin{lstlisting}[boxpos=t,aboveskip=0pt]
%module MyModule:
%  input x: {a: UInt, b: UInt}
%  output y:{a: UInt, b: UInt}
%  y <= x
%\end{lstlisting} &
%\raisebox{-6ex}{\quad into\quad}
%&
%&
%%[basicstyle=\footnotesize\ttfamily]
%\begin{lstlisting}[boxpos=t,aboveskip=0pt]
%module MyModule:
%  input x: {a: UInt, b: UInt}
%  output y:{a: UInt, b: UInt}
%  y.a <= x.a
%  y.b <= x.b
%\end{lstlisting}
%\end{tabular}
%\end{figure}

\paragraph{ExpandWhens} The ExpandWhens pass mainly replaces \prettyll{when} statements;
this is necessary because a physical circuit cannot contain conditional elements.
It also applies the so-called last-connect semantics.

For every ground-type writable (sub)component, ExpandWhens retains the last connect or \prettyll{is} \prettyll{invalid} statement
that appears under any combination of \prettyll{when} conditions in the module.
If the (sub)component is declared under multiple conditions,
ExpandWhens selects between them using a multiplexer.

\begin{figure}[h]
\centering
\begin{subfigure}[b]{0.4\textwidth}
         \centering
%[basicstyle=\footnotesize\ttfamily]
\begin{lstlisting}[boxpos=t,aboveskip=0pt,escapeinside={(*}{*)}]
module TestCondition:
  input in: UInt<32>
  input condition: UInt<1>
  output posv: UInt<32>
  output negv: SInt<33>
  posv <= 0 (* \label{lstinline-connect-before-when} *)
  when condition:
    posv <= in (* \label{lstinline-connect-in-when} *)
    negv is invalid (* \label{lstline-is-invalid} *)
  else:
    wire a: SInt<33> (* \label{lstline-wire-under-when} *)
    a <= neg in
    negv <= a
\end{lstlisting}
\vspace*{0.92\baselineskip}
\caption{}
\end{subfigure}
\hspace{4mm}
%     \hfill
\begin{subfigure}[b]{0.4\textwidth}
         \centering
%[basicstyle=\footnotesize\ttfamily]
\begin{lstlisting}[boxpos=t,aboveskip=0pt]
module TestCondition:
  input in: UInt<32>
  input condition: UInt<1>
  output posv: UInt<32>
  output negv: SInt<33>
  wire a: SInt<33>
  posv <= mux(condition, in, 0)
  a <= neg in
  negv <= validif(not condition, a)
\end{lstlisting}
\caption{}
\end{subfigure}
\caption{Examples for illustrating the ExpandWhens pass}
\label{fig-lofirrtl-exmp3}
\end{figure}
%
%\begin{figure}[h]
%\centering
%\begin{tabular}{cccc}
%%[basicstyle=\footnotesize\ttfamily]
%\begin{lstlisting}[boxpos=t,aboveskip=0pt,escapeinside={(*}{*)}]
%module TestCondition:
%  input in: UInt<32>
%  input condition: UInt<1>
%  output posv: UInt<32>
%  output negv: SInt<33>
%  posv <= 0 (* \label{lstinline-connect-before-when} *)
%  when condition:
%    posv <= in (* \label{lstinline-connect-in-when} *)
%    negv is invalid (* \label{lstline-is-invalid} *)
%  else:
%    wire a: SInt<33> (* \label{lstline-wire-under-when} *)
%    a <= neg in
%    negv <= a
%\end{lstlisting}&
%\raisebox{-14ex}{\qquad into\qquad}
%&
%&
%%[basicstyle=\footnotesize\ttfamily]
%\begin{lstlisting}[boxpos=t,aboveskip=0pt]
%module TestCondition:
%  input in: UInt<32>
%  input condition: UInt<1>
%  output posv: UInt<32>
%  output negv: SInt<33>
%  wire a: SInt<33>
%  posv <= mux(condition, in, 0)
%  a <= neg in
%  negv <= validif(not condition, a)
%\end{lstlisting}
%\end{tabular}
%\end{figure}

%\noindent

For instance, the ExpandWhens pass transforms the program shown in Figure~\ref{fig-lofirrtl-exmp3}(a)
into the one shown in Figure~\ref{fig-lofirrtl-exmp3}(b).
Note that according to the last connect semantics,
the connect statement in Line~\ref{lstinline-connect-before-when} of Figure~\ref{fig-lofirrtl-exmp3}(a)
is overridden by the one in Line~\ref{lstinline-connect-in-when} of Figure~\ref{fig-lofirrtl-exmp3}(a)
if \prettyll{condition} is true.
Because of this overriding,
ExpandWhens can only run after the connect statements have been expanded by ExpandConnects.

If, under some combination of conditions, the (sub)component is declared but is not connected to,
the compiler assumes that the programmer has forgotten a connect statement
and produces an error.
This error can be suppressed by an \prettyll{is invalid} statement
(Line~\ref{lstline-is-invalid} of Figure~\ref{fig-lofirrtl-exmp3}(a)),
which explicitly indicates that the connect is missing on purpose,
e.g. because the result is not used or is irrelevant under some conditions.
The multiplexer is replaced by a \prettyll{validif} expression.
%, which may allow later optimization passes to remove superfluous calculations.

If a component is only declared under some condition
(e.g. \prettyll{wire a} in Line~\ref{lstline-wire-under-when} of Figure~\ref{fig-lofirrtl-exmp3}(a))),
it is nevertheless included in the LoFIRRTL result,
and the condition is ignored; no multiplexer is inserted upon connections.

%\paragraph{LowerTypes} The LowerTypes pass flattens bundles and vectors, after the connect and partial connect statements have been expanded by
%the ExpandConnects pass. Nevertheless, LowerTypes does not change the implicit bundles, e.g. module instances. After the application of the LowerTypes pass, a module instance is still a bundle of depth 1 (instance -> port).

%the type of a memory is still a bundle with depth 2 (mem -> port -> field), see MemPortUtils.memType


%\hide{
%\paragraph{RemoveReset} The RemoveReset pass removes the synchronous resets. If the register is driven by an asynchronous reset, it will be left unchanged. If it is a synchronous reset, the reset operation defined in the register and the connections to the register will be merged into multiplexer operations. For example, the register \begin{lstinline}!myreg!\end{lstinline} with synchronous reset signal \begin{lstinline}!rst!\end{lstinline} in the program
%\begin{lstlisting}
%  module MyModule :
%    input clk : Clock
%    input rst : UInt<1>
%    input in : UInt<8>
%    input init : UInt<8>
%    reg myreg : UInt<8>, clk with: (reset => (rst, init))
%    myreg <= in
%\end{lstlisting}
%is transformed into
%\begin{lstlisting}
%  module MyModule :
%    input clk : Clock
%    input in : UInt<8>
%    input init : UInt<8>
%    reg myreg : UInt<8>, clk
%    myreg <= mux(``h0'', init, in)
%\end{lstlisting}
%}

